% ============================================================
% 3.2.3.3 Inflation and Deflation — Worksheet
% AQA AS Economics
% Sources: output/authoritative/as_syllabus.json (3.2.3.3)
%          output/authoritative/sow.json (as > 3.2 > 3.2.3)
% Total marks: 60  |  Suggested time: 2 hours (inc. group discussion)
% ============================================================

\documentclass[12pt,a4paper]{article}
\usepackage[margin=1in,headheight=15pt]{geometry}
\usepackage{enumitem}
\usepackage{amsmath}
\usepackage{fancyhdr}
\usepackage{booktabs}
\usepackage{array}
\usepackage{tikz}
\usepackage{pgfplots}
\pgfplotsset{compat=1.18}
\RequirePackage{tcolorbox}
\tcbuselibrary{skins}

% Key Point — yellow; for the single most important idea on a slide
\newtcolorbox{keypoint}{
    enhanced,
    colback=yellow!12,
    colframe=yellow!60!black,
    fonttitle=\bfseries,
    title=Key Point,
    arc=2pt,
    boxrule=0.8pt,
}

% Formula — blue; for named equations and their variable key
\newtcolorbox{formula}[1][Formula]{
    enhanced,
    colback=blue!5,
    colframe=blue!45!black,
    fonttitle=\bfseries,
    title=#1,
    arc=2pt,
    boxrule=0.8pt,
}

% Example Question — teal; for exam-style questions posed to students
\newtcolorbox{examplequestion}[1][Example Question]{
    enhanced,
    colback=teal!6,
    colframe=teal!55!black,
    fonttitle=\bfseries,
    title=#1,
    arc=2pt,
    boxrule=0.8pt,
}

% Example Solution — green; always follows an examplequestion block
\newtcolorbox{examplesolution}[1][Solution]{
    enhanced,
    colback=green!6,
    colframe=green!45!black,
    fonttitle=\bfseries,
    title=#1,
    arc=2pt,
    boxrule=0.8pt,
}

% ---- Worksheet environments --------------------------------

% Scenario — blue; for group discussion prompts on worksheets
% Usage: \begin{scenario}{Scenario 1: Title with commas, fine} ... \end{scenario}
\newtcolorbox{scenario}[1]{
    enhanced,
    colback=blue!5,
    colframe=blue!45!black,
    fonttitle=\bfseries,
    title={#1},
    arc=2pt,
    boxrule=0.8pt,
}

% Instructions — blue; for worksheet-level instruction blocks
\newtcolorbox{instructions}{
    enhanced,
    colback=blue!5,
    colframe=blue!45!black,
    fonttitle=\bfseries,
    title=Instructions,
    arc=2pt,
    boxrule=0.8pt,
}

% Extract — grey; for data response source material
% Usage: \begin{extract}{Source: Author, Year} ... \end{extract}
\newtcolorbox{extract}[1]{
    enhanced,
    colback=gray!8,
    colframe=gray!50,
    fonttitle=\bfseries,
    title={#1},
    arc=2pt,
    boxrule=0.8pt,
}


\pagestyle{fancy}
\fancyhf{}
\lhead{Dr Jac Thomas}
\rhead{Topic index: 3.2.3.3}
\rfoot{Page \thepage}

\begin{document}

% ---- Title block -------------------------------------------
\begin{center}
    \Large\textbf{Inflation and Deflation}\\[5pt]
    \normalsize Topic index: 3.2.3.3 \quad|\quad AQA AS Economics\\[3pt]
    \textit{Total Marks: 60 \quad Suggested Time: 2 hours (including group discussion)}
\end{center}

\vspace{0.8em}

\begin{instructions}
    Answer all questions in the spaces provided. Show your working clearly for calculation questions. For the group discussion, record your group's agreed conclusions in the spaces given.
\end{instructions}

\vspace{1em}

% ============================================================
% SECTION A: KEY TERMS
% ============================================================
\section*{Section A: Key Terms (10 marks)}

\textit{Define each of the following terms. Each definition is worth 2 marks.}

\vspace{0.5em}

\begin{enumerate}[leftmargin=*, label=\textbf{\arabic*.}]

    \item \textbf{Inflation} \hfill (2 marks)\\[3cm]

    \item \textbf{Deflation} \hfill (2 marks)\\[3cm]

    \item \textbf{Disinflation} \hfill (2 marks)\\[3cm]

    \item \textbf{Demand-pull inflation} \hfill (2 marks)\\[3cm]

    \item \textbf{Cost-push inflation} \hfill (2 marks)\\[3cm]

\end{enumerate}

\newpage

% ============================================================
% SECTION B: GROUP DISCUSSION
% ============================================================
\section*{Section B: Group Discussion (approximately 20 minutes)}

\textit{In your groups, read each scenario. For each, decide:}
\begin{enumerate}[label=(\alph*), topsep=2pt]
    \item What type of inflation (or price change) does this represent?
    \item What is the main cause?
    \item What policy might the government or Bank of England use to address it?
\end{enumerate}

\vspace{0.8em}

\begin{scenario}{Scenario 1: The UK Energy Crisis (2021--22)}
    Following supply disruptions in global gas markets, UK household energy bills more than doubled between 2021 and 2022. At its peak in October 2022, UK CPI inflation reached 11.1\% --- the highest rate in over 40 years. Much of the rise was driven by surging energy and food prices.
\end{scenario}

\vspace{0.5em}
\textbf{Your analysis:}\\
Type of inflation: \dotfill \\[0.4em]
Main cause: \dotfill \\[0.4em]
Appropriate policy response: \dotfill \\[0.4em]
Justification: \dotfill

\vspace{1em}

\begin{scenario}{Scenario 2: Post-COVID Demand Surge (2021)}
    When COVID-19 restrictions were lifted in 2021, UK consumers rushed to spend the savings accumulated during lockdown. Demand surged for restaurants, travel, and goods. At the same time, supply chains were severely disrupted, meaning businesses could not easily increase output to meet demand. Inflation began to rise sharply.
\end{scenario}

\vspace{0.5em}
\textbf{Your analysis:}\\
Type of inflation: \dotfill \\[0.4em]
Main cause: \dotfill \\[0.4em]
Appropriate policy response: \dotfill \\[0.4em]
Justification: \dotfill

\vspace{1em}

\begin{scenario}{Scenario 3: Japan's Lost Decade (1990s--2000s)}
    From the mid-1990s, Japan experienced persistent deflation --- prices fell year after year. Rather than stimulating spending, this led consumers to delay purchases, waiting for prices to fall further. Firms cut production and investment. Even very low interest rates failed to kick-start demand for over a decade.
\end{scenario}

\vspace{0.5em}
\textbf{Your analysis:}\\
What does this scenario illustrate? \dotfill \\[0.4em]
Why might falling prices be harmful to an economy? \dotfill \\[0.4em]
Is this deflation or disinflation? Explain the difference: \dotfill \\[0.4em]
What might a government do to escape a deflationary spiral? \dotfill

\newpage

% ============================================================
% SECTION C: MULTIPLE CHOICE
% ============================================================
\section*{Section C: Multiple Choice (10 marks)}

\textit{Circle the correct answer. Each question is worth 1 mark.}

\vspace{0.5em}

\begin{enumerate}[leftmargin=*, label=\textbf{\arabic*.}]

    \item Which of the following correctly defines disinflation?
    \begin{enumerate}[label=(\Alph*), topsep=2pt]
        \item The price level is falling
        \item The rate of inflation is falling while prices are still rising
        \item The price level is stable
        \item Real incomes are rising
    \end{enumerate}

    \vspace{0.4em}

    \item Which scenario is most consistent with demand-pull inflation?
    \begin{enumerate}[label=(\Alph*), topsep=2pt]
        \item Global oil prices rise sharply following OPEC production cuts
        \item Consumer spending surges as the economy approaches full employment
        \item A sharp depreciation of sterling raises import costs
        \item Wages increase across all sectors of the economy
    \end{enumerate}

    \vspace{0.4em}

    \item On an AD/AS diagram, cost-push inflation is most accurately illustrated by:
    \begin{enumerate}[label=(\Alph*), topsep=2pt]
        \item A rightward shift of the AD curve
        \item A rightward shift of the SRAS curve
        \item A leftward shift of the SRAS curve
        \item A rightward shift of the LRAS curve
    \end{enumerate}

    \vspace{0.4em}

    \item The Consumer Prices Index (CPI) measures:
    \begin{enumerate}[label=(\Alph*), topsep=2pt]
        \item The price of goods produced by UK manufacturers
        \item Changes in the weighted average price of a basket of goods and services
        \item The total value of UK exports
        \item Changes in the money supply
    \end{enumerate}

    \vspace{0.4em}

    \item A rise in world oil prices would most directly affect UK inflation through:
    \begin{enumerate}[label=(\Alph*), topsep=2pt]
        \item Raising production costs, shifting SRAS leftward
        \item Increasing aggregate demand, shifting AD rightward
        \item Causing the pound sterling to appreciate
        \item Reducing the UK money supply
    \end{enumerate}

    \vspace{0.4em}

    \item Which of the following is best described as a cost-push factor?
    \begin{enumerate}[label=(\Alph*), topsep=2pt]
        \item A cut in income tax stimulating consumer spending
        \item A fall in interest rates encouraging household borrowing
        \item A rise in firms' wage costs following a large increase in the National Living Wage
        \item A rise in government infrastructure spending
    \end{enumerate}

    \vspace{0.4em}

    \item If inflation falls from 6\% to 2\%, this is an example of:
    \begin{enumerate}[label=(\Alph*), topsep=2pt]
        \item Deflation
        \item Disinflation
        \item Stagflation
        \item Reflation
    \end{enumerate}

    \vspace{0.4em}

    \item ``Deflationary policies'' are best described as policies that aim to:
    \begin{enumerate}[label=(\Alph*), topsep=2pt]
        \item Cause the price level to fall
        \item Eliminate unemployment
        \item Reduce the growth of aggregate demand in order to reduce inflation
        \item Increase the productive capacity of the economy
    \end{enumerate}

    \vspace{0.4em}

    \item Strong economic growth in a major UK trading partner is most likely to affect UK inflation by:
    \begin{enumerate}[label=(\Alph*), topsep=2pt]
        \item Reducing demand for UK exports and lowering AD
        \item Increasing demand for UK exports, raising AD and putting upward pressure on prices
        \item Causing world commodity prices to fall
        \item Causing the pound to depreciate
    \end{enumerate}

    \vspace{0.4em}

    \item An economy experiencing simultaneously rising prices and falling output is most likely suffering from:
    \begin{enumerate}[label=(\Alph*), topsep=2pt]
        \item Demand-pull inflation caused by rising consumer spending
        \item Cost-push inflation from rising input costs shifting SRAS leftward
        \item A positive output gap driven by expansionary policy
        \item Deflation caused by falling aggregate demand
    \end{enumerate}

\end{enumerate}

\newpage

% ============================================================
% SECTION D: QUANTITATIVE SKILLS
% ============================================================
\section*{Section D: Quantitative Skills (8 marks)}

\textit{Show all working clearly. State the formula before substituting values.}

\vspace{0.8em}

The table below shows price data for a simplified price index for Econland, using 2022 as the base year. The index uses a weighted basket of four categories.

\vspace{0.5em}

\begin{table}[h]
    \centering
    \renewcommand{\arraystretch}{1.3}
    \begin{tabular}{@{}lrrrr@{}}
        \toprule
        \textbf{Category} & \textbf{Weight (\%)} & \textbf{Price (2022)} & \textbf{Price (2023)} & \textbf{Price Index (2023)} \\ \midrule
        Food        & 25 & £4.00  & £4.80  & ?   \\
        Housing     & 30 & £800   & £880   & 110 \\
        Transport   & 25 & £60.00 & £64.80 & ?   \\
        Clothing    & 20 & £50.00 & £52.50 & 105 \\ \bottomrule
    \end{tabular}
\end{table}

\vspace{0.8em}

\begin{enumerate}[leftmargin=*, label=\textbf{\arabic*.}]

    \item Calculate the price index for food in 2023, using 2022 as the base year (= 100). \hfill (2 marks)\\[0.2em]
    \textit{Show your working:}
    \vspace{3cm}

    \item Calculate the price index for transport in 2023. \hfill (1 mark)\\[0.2em]
    \textit{Show your working:}
    \vspace{2cm}

    \item Using all four price indices (your answers to questions 1 and 2, plus the two provided), calculate the overall price index (CPI) for Econland in 2023. \hfill (2 marks)\\[0.2em]
    \textit{Show your working:}
    \vspace{3.5cm}

    \item Calculate the rate of inflation in Econland between 2022 and 2023. \hfill (1 mark)\\[0.2em]
    \textit{Show your working:}
    \vspace{2cm}

    \item Between 2023 and 2024, inflation in Econland fell to 3\%. Explain whether this represents inflation, deflation, or disinflation, and justify your answer. \hfill (2 marks)\\[4cm]

\end{enumerate}

\newpage

% ============================================================
% SECTION E: DIAGRAMS
% ============================================================
\section*{Section E: Diagrams (10 marks)}

\begin{enumerate}[leftmargin=*, label=\textbf{\arabic*.}]

    \item Draw a fully labelled AD/AS diagram to illustrate demand-pull inflation. Your diagram should show the initial equilibrium, the shift that causes inflation, and the new equilibrium at a higher price level. \hfill (5 marks)

    \vspace{0.4em}
    \begin{center}
        \begin{tikzpicture}
            \begin{axis}[
                axis lines = middle,
                xlabel = {Real GDP},
                ylabel = {Price Level},
                xmin=0, xmax=10, ymin=0, ymax=10,
                xtick=\empty, ytick=\empty,
                width=10.5cm, height=7cm,
                xlabel style={at={(axis description cs:1,0)}, anchor=north east, yshift=-10pt, font=\small},
                ylabel style={at={(axis description cs:0,1)}, anchor=south east, font=\small},
            ]
            % Student draws here
            \end{axis}
        \end{tikzpicture}
    \end{center}

    \vspace{0.6em}

    \item Draw a fully labelled AD/AS diagram to illustrate cost-push inflation. Your diagram should show the initial equilibrium, the shift in the relevant curve, and the new equilibrium. Indicate what happens to both the price level and real GDP. \hfill (5 marks)

    \vspace{0.4em}
    \begin{center}
        \begin{tikzpicture}
            \begin{axis}[
                axis lines = middle,
                xlabel = {Real GDP},
                ylabel = {Price Level},
                xmin=0, xmax=10, ymin=0, ymax=10,
                xtick=\empty, ytick=\empty,
                width=10.5cm, height=7cm,
                xlabel style={at={(axis description cs:1,0)}, anchor=north east, yshift=-10pt, font=\small},
                ylabel style={at={(axis description cs:0,1)}, anchor=south east, font=\small},
            ]
            % Student draws here
            \end{axis}
        \end{tikzpicture}
    \end{center}

\end{enumerate}

\newpage

% ============================================================
% SECTION F: SHORT RESPONSE
% ============================================================
\section*{Section F: Short Response (7 marks)}

\begin{enumerate}[leftmargin=*, label=\textbf{\arabic*.}]

    \item Explain the difference between deflation and disinflation. In your answer, refer to what happens to the price level in each case. \hfill (4 marks)\\[6.5cm]

    \item Explain one way in which a rise in world commodity prices could lead to higher inflation in the UK. \hfill (3 marks)\\[5.5cm]

\end{enumerate}

\newpage

% ============================================================
% SECTION G: ESSAY
% ============================================================
\section*{Section G: Essay (15 marks)}

\begin{enumerate}[leftmargin=*, label=\textbf{\arabic*.}]

    \item Evaluate the view that rising aggregate demand is the main cause of inflation in the UK economy. \hfill (15 marks)

    \vspace{0.8em}

    \begin{keypoint}
        \textbf{Evaluation tips:} Consider both demand-pull and cost-push causes of inflation. Think about when demand-pull inflation is most likely (near full employment/positive output gap) and when cost-push may dominate (supply shocks, commodity price spikes). Consider how global factors --- world commodity prices, the exchange rate, and economic conditions abroad --- can cause UK inflation independently of domestic demand. Can one type of cause be clearly identified as ``more important''? Does it depend on the state of the economy and the time period? Think about the 2008--09 vs 2021--23 inflation episodes and which cause seems more relevant in each case.
    \end{keypoint}

    \vspace{14cm}

\end{enumerate}

\vfill
\begin{center}
    \textbf{[Total: 60 marks]}
\end{center}

\end{document}
